\chapter{Ophthalmologist}
\index{Ophthalmologist}

\medskip
\noindent\textit{\textbf{Under review.} This chapter is an AI-assisted educational draft; it is not a substitute for professional medical advice.}

\section*{Definition and Overview}
a medical doctor trained to diagnose and treat eye disease and perform eye surgery. Individual assessment is important because symptoms and treatment vary with the cause.

\section*{Clinical Features}
Visual loss or distortion, eye pain, redness, flashes or floaters, double vision, or other eye symptoms depending on the condition.

\section*{When to Seek Medical Care}
Seek medical review for new, persistent, worsening, or function-limiting symptoms. Seek emergency help for severe breathing difficulty, collapse, sudden neurological change, severe chest or abdominal pain, or any rapidly worsening illness.

\section*{Diagnosis and Investigations}
Visual acuity, pupils, eye movements, slit-lamp and dilated retinal examination, intraocular pressure, and imaging or visual-field tests as indicated.

\section*{Management}
Treatment may include glasses or medicines, laser or injections, surgery, or referral to another service; sudden visual loss or severe eye pain needs urgent assessment.

\section*{Complications}
Possible complications depend on the cause and may include progression of the condition, effects on other organs, treatment adverse effects, and reduced quality of life. Early assessment can reduce preventable harm.

\section*{Prognosis and Outcomes}
Outcome depends on the underlying cause, severity, complications, and response to treatment. Discuss individual prognosis with the treating clinician.

\section*{Prevention and Self-Care}
Follow the treatment plan, keep relevant appointments, avoid known triggers, protect affected areas from injury, and seek help early if symptoms worsen. Do not start or stop prescription treatment without clinical advice.
\section*{Safety-netting}
Seek urgent review for uncontrolled bleeding, severe pain, fever, spreading redness, wound separation, new neurological symptoms, breathing difficulty, or collapse.
